\chapter{Setting Up The System}
\label{ch:system_setup}

\section{System Setup}

This section explains, in simple steps, how to prepare and start the \textbf {License Plate Detection System}. Follow the setup steps below to get the system ready for usage:

\vspace{0.4cm}
\noindent
\textbf{Step 1: Place and Connect the Hardware}
\begin{itemize}
	\item Place the \textbf{System} on a stable, well-ventilated surface.
	\item Mount the \textbf{camera} on its stand and point it toward the vehicle lane or entry gate.
	\item Connect the camera to the System using the USB port.
	\item Connect the \textbf{HDMI cable} from the system to the monitor.
	\item Insert the \textbf{microSD card} into the system.
\end{itemize}

\vspace{0.3cm}
\noindent
\textbf{Step 2: Start the System}
\begin{itemize}
	\item Plug in the \textbf{power adapter} and turn on the power switch. 
	\item Wait a few seconds for the system to start.  
	\item The display will show the live camera view once it is ready.
\end{itemize}

\vspace{0.3cm}
\noindent
\textbf{Step 3: Verify Operation}
\begin{itemize}
	\item The system will now be able to detect number plates.Any detected car license plates are logged with timestamps and saved.
	\item Adjust the camera height or tilt if the plates are not clearly visible.
\item Ensure good lighting at the vehicle gate area.
For accurate plate reading, maintain an illumination level between 
\textbf{250 lux (minimum)} and \textbf{900 lux (maximum)} — with an 
optimal range of \textbf{400–700 lux}. 
Avoid shining lights directly into the camera to prevent glare or overexposure.
\end{itemize}

\vspace{0.3cm}
\noindent
\textbf{Step 4: During and After Use}
\begin{itemize}
	\item The system runs continuously and does not need supervision.  
	\item After use, simply turn off the power supply.  
	\item All data is stored safely for later review or reporting.
\end{itemize}

\vspace{0.5cm}
\noindent

\vspace{0.8em}
\noindent
\section{Camera Placement and Setup:}
\begin{itemize}
	\item Mount the camera about \textbf{2.0–2.5 metres above ground level}, facing the vehicle lane at a \textbf{30–40° angle}.
	\item Keep a distance of about \textbf{3–5 metres} between the camera and the point where vehicles stop or slow down.
	\item Avoid placing the camera directly opposite bright lights or reflective surfaces.
	\item Ensure there is enough ambient light. During night hours, add a small \textbf{LED lamp or floodlight} near the gate for better clarity.
	\item Keep the camera lens clean and free of dust, water spots, or scratches.
	\item Do not position the camera behind glass or near strong sunlight to prevent reflections.
\end{itemize}

\vspace{0.8em}

\section*{Important Tips}

\begin{itemize}
	\item Place the device on a stable, flat surface with good ventilation.
	\item Make sure the \textbf{power supply (5 V / 4 A)} is connected firmly and the cable is not loose.
	\item Keep the \textbf{camera lens clean} for accurate plate reading.
	\item Do not block the openings on the aluminum case; it keeps the system cool.
	\item Avoid exposing the device to rain, moisture, or direct sunlight.
	\item Connect the HDMI cable to your monitor before switching on the power.
	\item Wait until the indicator light turns on — this means the system is ready.
	\item If the display stays blank, check power and camera connections first.
\end{itemize}

\vspace{0.5em}
\noindent
These simple steps help ensure your system starts safely and performs reliably every time.

\vspace{0.8cm}

\begin{tcolorbox}[
	colback=yellow!10,
	colframe=yellow!50!black,
	title={Important to Know Before Use},
	fonttitle=\bfseries,
	left=6pt,
	right=6pt,
	top=6pt,
	bottom=6pt,
	boxrule=0.8pt
	]
	\begin{itemize}
		\item This system records number plates for time tracking and billing at the airport gate.
		\item Do not move or touch the camera once it is aligned. Misalignment can cause unreadable plates.
		\item If a plate is too dirty or not readable, you may need to type it in manually later.
		\item Emergency and official vehicles (police, ambulance, etc.) are recognised as \textbf{no charge}.
		\item Turn off the system only after normal shutdown. Unplugging power while it is running can cause data loss.
	\end{itemize}
\end{tcolorbox}

%\vspace{0.6cm}
%\noindent
%\textit{For detailed instructions on setup and use, please refer to \autoref{ch:system_setup}.}
%
%\vspace{0.5em}
%\noindent
%These simple steps help ensure your system starts safely and performs reliably every time.
%
